% Vorlage für eine Bachelor- oder Masterthesis
%% Dipl.-Ing. Stefan Averkamp
%% Version 2.5
%% 18.03.2021
%% ----------------------------------------------------------------------------

\documentclass[
cmyk,
pdftex,
a4paper,           % DinA4 Papier
%12pt,             % Schriftgröße
oneside,           % oneside, twoside
openright,         % openleft, openright, openany
numbers=noenddot,  % keinen Punkt hnter Kapitel- oder Sectionnummern
chapterprefix=off, % Kapitel wird vor die Kapitelnummer geschrieben
]{scrbook}
%% Preambel einbinden -----------------------------------------------------------------------------
\UseRawInputEncoding
\setcounter{secnumdepth}{3}         % 3 levels in section numbering

%% Useful packages -------------------------------------------------------------------------------
\usepackage[english]{babel}         % English language
\usepackage{scrhack}
\usepackage{placeins}
\usepackage{subcaption}
\usepackage{lmodern}
\usepackage{cmbright}
\usepackage{enumitem}               % Package for customized lists
\usepackage[OT1]{fontenc}
\usepackage{blindtext}              % Use \blindtext to insert dummy text
\usepackage{mdwlist}                % More compact lists with itemize and co
\usepackage{graphicx}               % Graphics
\usepackage{epstopdf}               % Automatically convert eps images to pdf
\usepackage{ltxtable}               % Combines longtable and tabularx
\usepackage{multirow}               % Multi-row cells in tables
\usepackage[dvipsnames,table,xcdraw]{xcolor}  % Alternating colors in tables
\usepackage{rotating}               % Rotate text and other elements
\usepackage{multicol}               % Multicolumn formatting
\usepackage{amsmath}                % Mathematical formulas
\usepackage{amsthm}                 % Theorem style
\usepackage{mathtools}
\usepackage{amssymb}                % Mathematical symbols
\usepackage{wrapfig}                % Wrapping figures
\usepackage{float}                  % For [h!] or [H] positioning
\usepackage{booktabs}               % For tables from Excel to LaTeX
\usepackage{eurosym}                % Use Euro symbol
\usepackage{url}                    % URLs
\usepackage{etoolbox}
\appto\UrlBreaks{\do\a\do\b\do\c\do\d\do\e\do\f\do\g\do\h\do\i\do\j
    \do\k\do\l\do\m\do\n\do\o\do\p\do\q\do\r\do\s\do\t\do\u\do\v\do\w
    \do\x\do\y\do\z}
\usepackage{adjustbox}
\usepackage{fancybox}               % Boxes and frames
\usepackage{caption}                % Extended settings for captions
\usepackage{array}                  % Extended options for arrays and tables
\usepackage{icomma}                 % Reduces space after comma in decimal numbers
\usepackage{setspace}               % Define line spacing

%% PDF options (links and more) -----------------------------------------------------------------
\usepackage{pdfpages}
\usepackage[utf8]{inputenc}         % UTF-8 encoding
\usepackage[english=british]{csquotes}
\usepackage[pdftex, raiselinks, pdfpagelabels, hypertexnames=false, pdfstartview={Fit}]{hyperref}
\definecolor{lightblue}{rgb}{0.9,0.9,1}

%% PDF metadata
\hypersetup{
    pdfauthor = {Nico Bierbaum},
    pdftitle =  {Scientific Work},
    pdfsubject = {Density and Shrinkage Evaluation of 316L Printed Parts in the FDM Process},
    pdfkeywords = {3D printing, filament, metal, density, tensile test, shrinkage},
    pdfcreator = {LaTeX with hyperref package}
}

%% Header and footer -----------------------------------------------------------------------------
\usepackage[headsepline]{scrlayer-scrpage}
\pagestyle{scrheadings}
\clearpairofpagestyles
\ihead{\headmark} % Redefine
\ohead{\pagemark}

%% Page margins -----------------------------------------------------------------------------------
\usepackage[a4paper,inner=4cm,outer=2.5cm,top=2cm,bottom=2cm,includeheadfoot]{geometry}

%% Bibliography -----------------------------------------------------------------------------------
\usepackage[style=numeric-comp,sorting=none,backend=biber]{biblatex}

\DeclareFieldFormat{labelalpha}{\textsc{#1}}    % Sets the author in small caps 
\renewcommand*{\mkbibnamefamily}[1]{\textsc{#1}}  % Sets the last name in small caps 
\renewcommand*{\mkbibnamegiven}[1]{\textsc{#1}}  % Sets the first name in small caps 
\renewcommand*{\mkbibnameprefix}[1]{\textsc{#1}}  % Sets the prefix in small caps 
\DeclareFieldFormat{title}{``#1''}      % Sets the title in quotation marks
\DeclareFieldFormat{date}{({#1})}               % Sets the date in parentheses
\makeatletter
\def\blx@maxline{77}
\makeatother
\addbibresource{LiteratureDB.bib} % Load bibliography

%% Glossary and abbreviations ---------------------------------------------------------------------
\usepackage{nomencl}
% English title
\renewcommand{\nomname}{Glossary and List of Abbreviations}
% Dots between abbreviation and explanation
\setlength{\nomlabelwidth}{.20\hsize}
\renewcommand{\nomlabel}[1]{#1 \dotfill}
% Reduce line spacing
%\setlength{\nomitemsep}{-\parsep}
\makenomenclature
\setlength{\headsep}{1.5\baselineskip}

%% ---------------------------------------------------------------------------------------------------
%% Create glossary and abbreviations: ------------------------------------------------------------------
%% 1. Run LATEX
%% 2. Open a command prompt and navigate to the folder with the "thesis.tex" file.
%% 3. Run the command   makeindex thesis.nlo -s nomencl.ist -o thesis.nls
%% 4. Run LATEX again. The glossary and abbreviations should now be present.

%% Source code -------------------------------------------------------------------------------------
\usepackage{listings}
\lstset{ %
    basicstyle=\scriptsize,         % the size of the fonts that are used for the code
    numbers=left,                   % where to put the line-numbers
    numberstyle=\footnotesize,      % the size of the fonts that are used for the line-numbers
    backgroundcolor=\color{white},  % choose the background color. You must add \usepackage{color}
    commentstyle=\color{green},
    stringstyle =\color{orange},
    xleftmargin=1.5em,
    xrightmargin=1em,
    frame=tb,                       % adds a frame around the code
    tabsize=2,                      % sets default tabsize to 2 spaces
    captionpos=b,                   % sets the caption-position to top or bottom
    breaklines=true,                % sets automatic line breaking
    breakatwhitespace=false         % sets if automatic breaks should only happen at whitespace
    literate =%
     {#}{\#}1
}
\renewcommand*{\lstlistingname}{Source Code}
\renewcommand*{\lstlistlistingname}{List of Source Codes}

%% Tikz for (complex) colorful drawings ---------------------------------------------------------
\usepackage{tikz}
\usetikzlibrary{positioning,mindmap,shapes,shapes.multipart,shadows,arrows,patterns,topaths}
\usepackage{pgfplots}
\pgfplotsset{width=7cm,compat=newest}

\definecolor{yellow}{HTML}{FFFCCC}
\definecolor{green}{HTML}{CFFFCC}
\definecolor{blue}{HTML}{CCE9FF}
\definecolor{darkblue}{HTML}{99b7EE}
\definecolor{red}{HTML}{FFA8A8}
\definecolor{orange}{HTML}{FFD28F}

\lstset
{
    breaklines=true,
    tabsize=3,
    showstringspaces=false
}

\lstdefinestyle{Common}
{
    extendedchars=true,
    language={[Visual]Basic},
    frame=single,
    %===========================================================
    framesep=3pt,%expand outward.
    framerule=0.4pt,%expand outward.
    xleftmargin=3.4pt,%make the frame fits in the text area. 
    xrightmargin=3.4pt,%make the frame fits in the text area.
    %=========================================================== 
    rulecolor=\color{Red}
}

\lstdefinestyle{A}
{
    style=Common,
    backgroundcolor=\color{Yellow!10},
    basicstyle=\scriptsize\color{Black}\ttfamily,
    keywordstyle=\color{Orange},
    identifierstyle=\color{Cyan},
    stringstyle=\color{Red},
    commentstyle=\color{Green}
}

\lstdefinestyle{B}
{
    style=Common,
    backgroundcolor=\color{Black},
    basicstyle=\scriptsize\color{White}\ttfamily,
    keywordstyle=\color{Orange},
    identifierstyle=\color{Cyan},
    stringstyle=\color{Red},
    commentstyle=\color{Green}
}
% Default fixed font does not support bold face
\DeclareFixedFont{\ttb}{T1}{txtt}{bx}{n}{12} % for bold
\DeclareFixedFont{\ttm}{T1}{txtt}{m}{n}{12}  % for normal

% Custom colors
\usepackage{color}
\definecolor{deepblue}{rgb}{0,0,0.5}
\definecolor{deepred}{rgb}{0.6,0,0}
\definecolor{deepgreen}{rgb}{0,0.5,0}

% Python style for highlighting
\newcommand\pythonstyle{\lstset{
language=Python,
basicstyle=\ttm,
morekeywords={self},              % Add keywords here
keywordstyle=\ttb\color{deepblue},
emph={MyClass,__init__},          % Custom highlighting
emphstyle=\ttb\color{deepred},    % Custom highlighting style
stringstyle=\color{deepgreen},
frame=tb,                         % Any extra options here
showstringspaces=false
}}


% Python environment
\lstnewenvironment{python}[1][]
{
\pythonstyle
\lstset{#1}
}
{}

\begin{document}

%% Titelseite einbinden ---------------------------------------------------------------------------
\thispagestyle{empty}
\newgeometry{
    top=25mm, 
    bottom=7mm,
    left=46mm,
    right=36mm
}
%Logos------------------------------------------------------------------------------
%Logos have a defined width of 6cm, the box containing the logo is 7cm
\vspace*{-2cm} % may be useful to place logos higher
\setlength{\unitlength}{1cm}
\hfill
\begin{minipage}[t]{4cm} 
    \begin{center}
       \includegraphics[width=40mm]{bilder/Firmenlogo.png}
    \end{center}
\end{minipage}
 %% University Logo ----------------------------------------------------------------------
 \hfill  % This adds space and pushes the next minipage to the right
 \begin{minipage}[t]{7cm}
    \begin{flushright}
       \includegraphics[width=66mm]{bilder/fhm_Logo_CMYK_30mm.pdf}
    \end{flushright}
 \end{minipage}
%-----------------------------------------------------------------------------------

\begin{minipage}[t]{11cm}
\vspace{2cm}                    
\begin{Huge}
\begin{center}\textsf{\textbf{Master's Thesis}}\end{center}
\end{Huge}


\vspace{2cm}                                
\begin{Large}
\hfill 
\begin{center}

\begin{large}
    Nico Bierbaum
\end{large}

\vspace{2cm} %3cm
%\begin{minipage}[t][2.5cm]{14cm}
\begin{flushright}
    \hfill{Development of a Robot-Based Method for Characterizing Single-Use Bag Mixers}\\
\end{flushright}            
%\end{minipage}
\end{center}    
\end{Large}

\begin{large}
\vspace{30mm}  %3cm
%\begin{minipage}[t][1.5cm]{14cm}
\hfill Supervisor: Prof. Dr.-Ing. Stephan Behr

\hfill Marius Kling M.Sc.\\
%\end{minipage}
\end{large}

\vspace{35mm}
\hfill Leverkusen, \today
\end{minipage}
\vspace{20mm}

\begin{minipage}[t]{136.5mm}
\begin{flushright}
\includegraphics[width=82mm]{bilder/MB_CMYK.pdf}
\end{flushright}
\end{minipage}
\restoregeometry
%--------------------------------------------------------------------------------------------------------


%% Sperrvermerk, Eidesstattliche Erklärung, Vorwort -----------------------------------------------
\frontmatter

\pagestyle{scrheadings}
\pagenumbering{Roman}
%\input{inhalt/kapitel1.tex}

%\input{inhalt/Sperrvermerk.tex}				nicht notwenig???
%%Eidesstattliche Erklärung-------------------------------------------------------------------------------
\chapter*{Eidesstattliche Erklärung}
\addcontentsline{toc}{chapter}{Eidesstattliche Erklärung}

Ich erkläre hiermit an Eides statt, dass ich die vorliegende Hausarbeit selbstständig und ohne unzulässige fremde Hilfe angefertigt habe. Die verwendeten Literaturquellen sind im Literaturverzeichnis vollständig zitiert.





\vspace{6cm}
Steinfurt, \today
	nicht notwendig???
%% Vorwort -----------------------------------------------------------------------------------------
\chapter*{Vorwort}

Schreiben wir den Bums hier ein bisschen wie Werbung oder soll das komplett sachlich und nüchtern passieren?

\addcontentsline{toc}{chapter}{Vorwort}



%% Inhaltsverzeichnis einbinden -------------------------------------------------------------------
\pdfbookmark[0]{Inhaltsverzeichnis}{tableofcontents}
\tableofcontents

%% Verzeichnisse einbinden  -----------------------------------------------------------------------
\listoffigures              %% Abbildungsverzeichnis
\listoftables               %% Tabellenverzeichnis
%\lstlistoflistings          %% Quelltextverzeichnis
\nomenclature{AM}{Additive Manufacturing - additive Fertigung}



\cleardoublepage% or \clearpage
\markboth{\nomname}{\nomname}

\printnomenclature
\addcontentsline{toc}{chapter}{Glossar und Abkürzungsverzeichnis}
  %% Abkuerzungsverzeichnis
%% Eigentlicher Inhalt der Arbeit -----------------------------------------------------------------
%% Einbinden der einzelnen Kapitel. Jedes Kapitel in eine seperate Datei speichern. ---------------
\mainmatter

\chapter{Einleitung}
\label{einleitung}

\subsection*{Disclaimer zur Verwendung des generischen Maskulinums}
Für die Lesbarkeit und bessere Verständlichkeit dieses Businessplans haben wir uns dazu entschieden, das generische Maskulinum zu verwenden. Wir möchten darauf hinweisen, dass wir damit keinesfalls eine Diskriminierung oder Geringschätzung von Personen anderer Geschlechter beabsichtigen. Selbstverständlich sind alle im Text verwendeten Bezeichnungen geschlechtsneutral zu verstehen und gelten gleichermaßen für Frauen und Männer.

\pagebreak

\section{Hintergrund und Motivation}
Additive Fertigung \textit{engl. Additive Manufacturing (AM)} wurde in den 1990ern eingeführt und beschreibt eine bis dahin neuartige Technologie 3D-Objekte in einem Schicht-für-Schicht Verfahren zu fertigen. Im Gegensatz zur additiven Fertigung [AM] steht die subtraktive Fertigung, auch als konventionelle Fertigung bezeichnet. Bei der die gewünschten Bauteile durch Abtrag von Material (z.B. beim Fräsen oder Drehen) hergestellt werden. Durch die additive Fertigung ergibt sich eine hohe Flexibilität und Designfreiheit, dies ist insbesondere für Prototypenherstellung, als auch immer mehr in der Serienproduktion von Bedeutung \autocite{Prof.Dr.Ing.ChristianSeidel.2023}.
Zu Beginn ist diese Technologie hauptsächlich in der Prototypenphase eingesetzt, da sich die Materialienauswahl auf Polymere beschränkt. So gab es vor allem viele Kunststoffe zur Auswahl, aber keine Materialien, die einem Stahl gleichkommen.
Heute ist AM im weiteren Sinne auch für die Herstellung von funktionellen Teilen eingesetzt.  Firmen aus den Bereichen Luft- und Raumfahrt, Automobilherstellung, Energie und Medizin sind interessiert an dem Einsatz der AM zur Herstellung verschiedener Bauteile.
AM bringt folgende Vorteile mit sich:
\begin{itemize}
    \item Große Freiheit im Design der Bauteile, was eine individuelle Anpassung an die Produktion ermöglicht. Zudem ist die Erstellung komplexer Bauteile deutlich einfacher als mit konventionellen Methoden.
    \item Der Konstruktionsprozess und die Prototypenphase sind deutlich schneller. Einzelne Prozessschritte sind schneller.
\end{itemize}

\section{Zielsetzung der Arbeit}

Das Ziel dieser Wissenschaftlichen Arbeit besteht darin, die Dichteeigenschaften und das Schrumpfungsverhalten des 316L (Edelstahl) der Firma \textit{PT+A} zu evaluieren. 
Genauer werden die Maße und das Gewicht in allen drei Verarbeitungsstufen aufgenommen um daraus die Dichte zu bilden. Dadurch wird deutlich wieviel Material durch entbinden und sintern in Verlust gerät.
Im Idealfall ist die Dichte des gesinterten Bauteils ähnlich zu dem von handelsüblichem 316L.
Zusätzlich werden Zugproben gedruckt, die Aufschluss über die mechanischen Kennwerte des Materials geben. Diese werden auch mit handelsüblichem 316L verglichen.
    
\section{1.3 Aufbau der Arbeit}

Nach dieser Einleitung werden die technischen Grundlagen aufgearbeitet. Im Anschluss wird im dritten Kapitel wird der praktische Teil dieser Arbeit behandelt, dies umfasst den 3D-Druck und die optimierung der Druckparameter. Im vierten Kapitel werden die Ergebnisse der Dichtemessung und des Zugversuches diskutiert und verglichen. Zuletzt gibt es einen Ausblick.


\chapter{Theoretische Grundlagen}

Die additive Fertigung hat in den letzten Jahren Fortschritte gemacht und ist zu einem Bestandteil der modernen Fertigungstechnologie geworden. Innerhalb dieses Prozesses werden Bauteile über CAD \texttt{(computer aided design)} entworfen. Über Verfahren der additiven Fertigung werden diese Bauteile dann gefertigt. FDM-Druck hat sich als eine vielversprechende Methode etabliert, um komplexe dreidimensionale Objekte schichtweise aufzubauen. Während FDM-Druckverfahren traditionell auf Polymermaterialien ausgerichtet sind, hat die Weiterentwicklung dieser Technologie nun den Weg für den Einsatz von Metallfilamenten eröffnet. In diesem Kapitel werden die theoretischen Grundlagen erläutert.

\subsection{Additive Fertigung und der FDM-Prozess}

Additive Fertigung wird als Verfahren bezeichnet, bei dem Materialien entweder durch Verschmelzung, Bindung oder Verfestigung von Materialien wie flüssigen Harzen und Pulvern gemischt werden. 

\autocite{Osama2019}

2.2 Materialien für den FDM-Prozess\\
2.3 316L-Edelstahl als Druckmaterial\\
2.4 Dichte- und Schwindungsbegriffe\\
\chapter{Experimenteller Aufbau}

In diesem Kapitel wird darauf eingegangen, wie die Proben erstellt werden, welche Druckparameter eine große Rolle Spielen und wie diese gefunden werden. Ausserdem wird erklärt, welche Messverfahren und welche Messgeräte eingesetzt werden.\\


\section{Auswahl der Druckparameter}

Beim verwenden eines 3D-Druckers ist es notwendig die Prozessparameter passend zum gewählten Filament anzupassen. Besonders bei Filamenten die dem Standard (PLA, PETG, ABS, etc.) abweichen, ist dies eine Herausforderung, da ein voreingestelltes Profil für den gewünschten 3D-Drucker nicht garantiert ist. Im Vorfeld dieser Arbeit ist bereits eine Wissenschaftliche Arbeit zum Thema \glqq FDM-Druck mit PT+A 316L Metallfilament Untersuchung der Druckparameter am Raise 3D Pro 2\grqq erstellt. Das Ergebnis dieser Arbeit lässt sich in Tabelle aus \autoref{ursprüngliche Werte} darstellen.

\begin{table}[htbp]
    \centering
    \caption{Druckparameter - gegeben aus \autocite{M.Mickan}}
      \begin{tabular}{llr}
      \toprule
      \textbf{Slicing-Paramater} & \textbf{Empfehlung} & \multicolumn{1}{l|}{\textbf{Hinweise (IdeaMaker)}} \\
      \midrule
      Drucktemperatur & 130°C & \multicolumn{1}{p{12.555em}}{+- 5 °C möglich} \\
      Druckgeschwindigkeit & 30-60 mm/s & \multicolumn{1}{p{12.555em}}{Füllung schnell, Konturen langsam} \\
      Heizbetttemperatur & 40 °C & \multicolumn{1}{p{12.555em}}{Gute Ablösung und Schichthaftung} \\
      Rückzugsgeschwindigkeit & 20 mm/s & \multicolumn{1}{p{12.555em}}{0,5 mm Rückzugsmenge} \\
      Materialflussrate & \multicolumn{1}{r}{90\%} & \multicolumn{1}{p{12.555em}}{Bei Extrusionsbreite 0,4 mm} \\
      Füllflussrate & \multicolumn{1}{r}{90\%} & \multicolumn{1}{p{12.555em}}{Vorsicht: Reiter „Fortgeschritten“ überschreibt Flussrateneinstellungen} \\
      Xy-Größenkompensation für Konturen & 0,1 mm &  \\
      Xy-Größenkompensation für Bohrungen & 0,06 mm &  \\
      Füllüberlappung & \multicolumn{1}{r}{10\%} & \multicolumn{1}{p{12.555em}}{Höher falls Ghosting / Pillowing eintritt} \\
      Schichthöhe & 0,3 mm &  \\
      \bottomrule
      \end{tabular}%
    \label{ursprüngliche Werte}%
  \end{table}%
  
Essentielle Parameter werden in dieser Arbeit jeweils mit Testdrucken bestimmt. Zur Bestimmung der Drucktemperatur wird ein sogenannter \textit{Heat Tower} gedruckt. Dabei wird in bestimmten Abständen die Temperatur in 5°C-Schritten abgesenkt. Mit dem gedruckten Bauteil lässt sich dann die geeigneteste Temperatur ablesen. \Autocite{M.Mickan}.\\

Diese Werte werden somit zu Beginn dieser Arbeit übernommen und erste Druckversuche werden durchgeführt. Dazu wird ein Würfel mit einer Kantenlänge von 10mm gewählt. Das Ergebnis ist in \autoref{erster Druck - Isometrie} und \autoref{erster Druck - Seitenansicht} dargestellt.


\begin{figure}[h]
	\centering
	\includegraphics[width=0.5\linewidth]{bilder/Testdruck auf Raise Pro 3D.png}
        \caption[Testdruck von dem Raise3D Pro2 Plus - Isometrische Ansicht] {Testdruck von dem Raise3D Pro2 Plus - Isometrische Ansicht (Eigene Darstellung)}
	\label{erster Druck - Isometrie}
\end{figure}
\begin{figure}[h]
	\centering
	\includegraphics[width=\linewidth]{bilder/Testdruck auf Raise Pro 3D Seitenansicht - Draufsicht.png}
        \caption[Erster Testdruck mit dem Raise3D Pro2 Plus - Seitenansicht und Draufsicht] {Erster Testdruck mit dem Raise3D Pro2 Plus -Draufsicht (l) und Seitenansicht (r) (Geschwindigkeit: 60mm/s; Schichthöhe: 0,4mm) (Eigene Darstellung)}
	\label{erster Druck - Seitenansicht}
\end{figure}

Erkennbar ist hier deutlich, die ungenauigkeit des Drucks dargestellt. Dies lässt sich sehr gut in \autoref{erster Druck - Seitenansicht} an den Ecken in der rechten Draufsicht erkennen. Das Symptom der zu stark abgerundeten Ecken ist auf das Fehlen von \textit{linear advanced} aus \autoref{linear advanced} zurückzuführen. Die Extrusionsrate ist nicht an das Be- und Entschleunigungsverhalten des Druckkopfs angepasst. Dadurch ist an den Ecken, in denen die Geschwindigkeit langsamer ist, mehr Material aufgetragen. Die Einstellung des \textit{linear advanced} lässt sich bei diesem Drucker nicht einstellen und die \textit{Firmware} ist nicht \textit{Open Source}.

Bei Verringerung der Geschwindigkeit von 60mm/s auf 20mm/s ist dieser Effekt deutlich verringert. In \autoref{zweiter Druck - Isometrie} ist das Ergebnis der Verringerung der Geschwindigkeit und der Schichthöhe dargestellt. Doch diese Optimierung geht zu Lasten der gesamten Druckzeit. 

\begin{figure}[h]
	\centering
	\includegraphics[width=\linewidth]{bilder/2. Testdruck auf Raise Pro 3D Seitenansicht - Draufsicht.png}
        \caption[Zweiter Testdruck mit dem Raise3D Pro2 Plus - Isometrische Ansicht] {Zweiter Testdruck mit dem Raise3D Pro2 Plus -Draufsicht (l) und Seitenansicht (r) (Geschwindigkeit: 20mm/s; Schichthöhe: 0,1mm) (Eigene Darstellung)}
	\label{zweiter Druck - Isometrie}
\end{figure}

\subsection{Verwendung des Prusa i3 MK3S+}
\label{Drucker}

Mit dem Raise 3D Pro 2 sind die Ergebnisse aufgrund der fehlenden Einstellmöglichkeit des \textit{linear advanced} nicht zufriedenstellend. Da auch die verlängerte Druckzeit bei verringerter Geschwindigkeit nicht zufriedenstellend ist, wird in dieser Arbeit auf den \textit{i3 MK3S+} der Frima \textit{Prusa} zurückgegriffen. Dieser hat ebenfalls wie der \textit{Raise3D}-Drucker einen Druckkopf mit Direktextruder. Dies ist vom Filamenthersteller aufgrund des spröden Materials empfohlen.

\begin{figure}[h]
	\centering
	\includegraphics[width=\linewidth]{bilder/Erster Testdruck mit Prusa.png}
        \caption[Erster [20mm/s] und zweiter [50mm/s] Testdruck mit dem Prusa] {Erster [20mm/s] und zweiter [50mm/s] Testdruck mit dem Prusa (Eigene Darstellung)}
	\label{Erster Testdruck Prusa}
\end{figure}

Der erste Testdruck mit dem \textit{Prusa} ist in \autoref{Erster Testdruck Prusa} links dargestellt. Die Einstellungen gleichen den Einstellungen aus \autocite{M.Mickan}
(siehe \autoref{ursprüngliche Werte}). Die Geschwindigkeit ist jedoch auf 20mm/s reduziert und die Schichthöhe beträgt 0,1mm. Aufgrund des fehlerfreien Ergebnisses wird ebenso ein Würfel mit einer Geschwindigkeit von 50mm/s gedruckt. Dieser ist in \autoref{Erster Testdruck Prusa} rechts dargestellt. Somit stehen die Parameter für die Würfel und die liegenden Zugproben fest. Sie sind in \autoref{Würfelparameter} aufgefasst. 

\section{Probenherstellung und -messung}
Die Testwürfel sind in dem \textit{CAD}-System \textit{NX\textbf{VERSION??}} erzeugte Würfel mit einer Kantenlänge von 10mm. Diese sind dann als \textit{STL}-Datei exportiert und mit \textit{Ultimaker Cura Version 5.3.1} geslicet, also ein \textit{g-code} erzeugt. Wie bereits erwähnt wird für den Druck der \textit{Prusa MK3S+} mit einer \textbf{DÜSE??} verwendet. Auf dem originalen Druckbetts des \textit{Prusa} wird eine Schicht aus Flachkreppband der Marke \textit{Günter Seits \textbf{PRüfen!}} geklebt. Ohne dieses Flachkreppband haftet die erste Schicht zu gut an dem Druckbett, dadurch ist es unmöglich das Bauteil zerstörungsfrei von dem Druckbett zu lösen.\\
Die Würfel werden nun als Grünteil mit einem Digitalmessschieber der Marke \textit{Mitutoyo} gemessen. Die Proben werden einzeln mit einer \textit{Kern PRS 620-3} gewogen. Die Messwerte, samt errechntem Volumen und Dichte sind in \autoref{Grünteilmaße} aufgeführt.
\chapter{Ergebnisse und Diskussion}

Im Fokus dieses Kapitels stehen die Ergebnisse und Diskussionen im Kontext der Evaluierung von Dichteeigenschaften. Schrumpfungsverhalten und mechanische Kennwerte des 316L-Edelstahls von \textit{PT+A} für den 3D-Druck mittels FDM. Im vorangegangenen Kapitel wird die Erstellung der Proben und die wissenschaftlichen Grundlagen erläutert.

\section{Dichte- und Schwindungsauswertungen der gedruckten Proben}
\label{Dichte- und Schwindung}
In diesem Unterkapitel werden die gedruckten Würfel- und Zugproben hinsichtlich der Dichte untersucht. Diese Messung wird zuerst im Grün-, sowie im Braun-, als auch im späteren Sinterteil durchgeführt, damit die Schwindungen und die Massenverluste bestimmt werden kann. Die Ergebnisse dieser Auswertung ist in \autoref{Grünteilmaße} bis \autoref{Sinterteilmaße} dargestellt. Die Durchschnittswerte sind in \autoref{tab:Rohdaten} dargestellt. Sie geben einen Überblick über die das Verhalten der Würfel in jedem der drei Zustände.
Aus diesen Werten lassen sich die Schwindungen der einzelnen Parameter nach jedem Prozess ableiten. Dargestellt sind diese Werte in \autoref{tab:Schrumpfung}. Dargestellt sind hier die relative Schrumpfungen/Zunahmen in Vergleich zum Grünteil.\\
Die Ergebnisse der Dichtemessung ist in \autoref{DichtewerteTestwürfel} visuell dargestellt. Im Durchschnitt beträgt die Dichte 6,87 \(\text{g/cm}^3\). Nachfolgend werden diese Werte mit dem Datenblatt und anderen Studien verglichen. Es wird ebenfalls erklärt, wie das Verhalten zustande kommt.

\begin{figure}[h] 
  \centering
  \includegraphics[width=0.8\linewidth]{bilder/img_Dichtewerte der Testwürfel.png}
        \caption[Dichtewerte der Testwürfel als Metallteil] {Dichtewerte der Testwürfel als Metallteil (Quelle: eigene Darstellung)}
  \label{DichtewerteTestwürfel}
\end{figure}
%\FloatBarrier


\begin{table}[h]
  \centering
  \caption{Datentabelle der Würfelproben}
  \captionsetup{font=small} % Schriftgröße auf small setzen
  \footnotesize
  \begin{tabular}{ccccccc}
    \toprule
    \textbf{Teil} & \multicolumn{1}{c}{\textbf{Z [mm]}} & \multicolumn{1}{c}{\textbf{X [mm]}} & \multicolumn{1}{c}{\textbf{Y [mm]}} & \multicolumn{1}{c}{\textbf{Volumen [mm\(^3\)]}} & \multicolumn{1}{c}{\textbf{Masse [g]}} & \multicolumn{1}{c}{\textbf{Dichte [g/cm\(^3\)]}} \\
    \textbf{Grünteil} & 9,99 & 10,17 & 10,19 & 1035,89 & 4,75 & 4,59 \\
    \textbf{Braunteil} & 10,00 & 9,88 & 9,90 & 977,29 & 4,52 & 4,62 \\
    \textbf{Sinterteil} & 8,62 & 8,60 & 8,58 & 636,20 & 4,37 & 6,87 \\
    \bottomrule
  \end{tabular}%
  \label{tab:Rohdaten}%
\end{table}%
\FloatBarrier

\begin{table}[h]
    \centering
    \caption{Schrumpfung der Würfelproben im Braun- und Sinterteil}
      \begin{tabular}{ccccccc}
      \toprule
      \textbf{Teil} & \multicolumn{1}{c}{\textbf{Z [\%]}} & \multicolumn{1}{c}{\textbf{X [\%]}} & \multicolumn{1}{c}{\textbf{Y [\%]}} & \multicolumn{1}{c}{\textbf{Volumen [\%]}} & \multicolumn{1}{c}{\textbf{Masse [\%]}} & \multicolumn{1}{c}{\textbf{Dichte [\%]}} \\
        \textbf{Entbindern} & 0 & -2,9 & -2,9 & -5,8 & -4,8 & 0,8 \\
        \textbf{Sintern} & -13,7 & -15,4 & -15,8 & -38,6 & -8 & 50 \\
      \bottomrule
      \end{tabular}%
    \label{tab:Schrumpfung}%
  \end{table}%
  \FloatBarrier

Wie in \Autocite{Industrie4} und in \autoref{fig:Verlust} dargestellt, ist die Volumenschrumpfung nach dem Entbindeprozess deutlich geringer, als im Sinterprozess. Grund hierfür ist der Binder, der durch das Lösemittel gelöst wird. Dadurch entstehen verhältnismäßig kleine Volumenschwindungen. Die Metallpartikel liegen wenig verändert in dem selben Gefüge, das erklärt ein Volumeverlust von 5,8 \%. Nach dem Sintern rücken die Metallpartikel deutlich dichter zueinander, was zu einer Volumenschrumpfung von 38,6 \% führt. Der Verlust der Masse beträgt im Entbindeprozess 4,8 \% und im Sinterprozess 8 \%. Bei der Veränderung der Masse sieht es anders aus. Die Erklärung liegt in der Vorgehensweise des Entbindern. Das Lösemittel diffundiert in das Grünteil hinein und löst das Bindemittel. Dabei wird ein großer Teil abgetragen, dieser liegt aber hauptsächlich in der Randzone der Würfel (siehe \autoref{fig:Verlust}). In dieser Arbeit beläuft sich der Abtrag des Bindemittels aus dem Entbindeprozess knapp über die Hälfte des gesamten Massenverlustes.Nach dem Entbindern bleibt ein Gerüst aus Bindemittel, welches die Geometrie für das Sintern zusammenhält. Dies nennt man auch Rückgrat\textit{engl. Backbone} \Autocite{Industrie4}.
Im Sinterprozess wird dieser Rest an Binder entfernt. Dies führt dann zu einem Gewichtsverlust von 8 \%. \textbf{VERGLEICH?}\\
Die Dichte nimmt über beide Prozessschritte zu. Somit ist sie der einzige Parameter, der nicht abschrumpft, sondern zunimmt. Diese Entwicklung lässt sich darauf zurückführen, dass das Volumen deutlich stärker schrumpft als die Masse. Die Dichte, die durch die Division des Volumens durch die Masse berechnet wird, erfährt daher eine Steigerung.

\begin{figure}[h] 
  \centering
  \includegraphics[width=0.8\linewidth]{bilder/Kukla_Bild1.png}
        \caption[Darstellung Formgebung, Entbindern und Sintern] {Darstellung Formgebung, Entbindern und Sintern \autocite{Industrie4}}
  \label{fig:Verlust}
\end{figure}

Von den liegenden Zugproben wurden auch Masse und Gewicht gemessen. Daher lassen sich auch hier das Schrumpfungsverhalten darstellen. Diese Werte sind in \autoref{Schrumpfung Zugproben} angegeben. Dabei decken sich die Schrumpfung des Volumens mit denen der Würfel. Lediglich die Schrumpfung der Masse ist mit 12 \% genau 4 Prozentpunkte mehr als bei den Würfeln. Dies ist ein Zeichen für eine effektivere Entbinderphase. \textbf{Dichte der Zugproben ermitteln und einf.}


\begin{table}[h]
    \centering
    \caption{Schrumpfung der liegenden Zugproben im Sinterteil}
      \begin{tabular}{ccccccc}
      \toprule
      \textbf{Prozess} & \multicolumn{1}{c}{\textbf{Volumen [\%]}} & \multicolumn{1}{c}{\textbf{Masse [\%]}} & \multicolumn{1}{c}{\textbf{Dichte [\%]}} \\
        \textbf{Sinterteil} & -38,4 & -12 & +42,9 \\
      \bottomrule
      \end{tabular}%
    \label{Schrumpfung Zugproben}%
  \end{table}%
  \FloatBarrier

\section{Ergebnisse der Zugversuche}

In diesem Abschnitt werden die mechanischen Kennwerte des 316L-Edelstahls mittels eines Zugversuchs ermittelt. Dazu wird mithilfe des, in \autoref{Drucker} vorgestellten, FDM-Drucker Zugproben hergestellt. Diese Zugproben sind nach DIN 50125-E3x8x30 hergestellt.
Ziel dieser Arbeit ist ein Vergleich zwischen zwei verschiedene Ausrichtungen dieser Proben.
\begin{figure}[h] 
  \centering
  \includegraphics[width=0.4\linewidth]{bilder/Screenshot 2023-11-01 171423.png}
        \caption[Stehende Zugprobe im Cura-Slicer] {Stehende Zugprobe im Cura-Slicer (Quelle: \autocite{Prusa})}
  \label{stehende Zugprobe}
\end{figure}
\FloatBarrier

\textbf{WARUM VERGLEICH STEHENDE LIEGENDE ZUGPROBEN}
Die Orientierung der Werkstücke in der Slicing-Software spielt eine große Rolle. Um diese Unterschiede sichtbar zu machen, wird in dieser Arbeit der Vergleich zwischen \glqq liegenden\grqq und \glqq stehenden\grqq Zugproben angestellt. Die \glqq liegenden\grqq Zugproben liegen flach auf dem Druckbett, was zu großen Druckflächen pro Druckschicht führt (siehe \autoref{fig:liegendeZugproben}). Verglichen mit dieser Fertigungsmethode werden \glqq stehende\grqq Zugproben. In diesem Fall werden die Zugproben auf ihrer kleinsten Fläche aufgestellt (siehe \autoref{fig:stehendeZugproben}). Mit dieser Methode werden kleinere Flächeninhalte als bei den \glqq liegenden\grqq Zugproben erreicht.

Beim Zugversuch werden jeweils die dickeren Enden der Zugprobe in eine Vorrichtung eingespannt. Die Zugmaschine bringt eine Kraft an beide Enden auf, die dann eine Längung verursacht, die dann bis zum Bruch führt.\\
Von Interesse ist in dieser Arbeit der Unterschied zwischen eine Kraft die längs der Druckschicht angebracht ist (siehe \autoref{fig:Kraftlängs}), zu einer Kraft die quer an die Druckschicht angebracht wird (siehe \autoref{fig:KraftQuer}). \textbf{Unzufriedener Satzbau}
Dieser Vergleich bringt Aufschluss über Optimierungsmöglichkeiten, die über Orientierung der Werkstücke erreicht werden kann. Ausserdem werden die erreichten Werte in \autoref{sec:Vergleich} mit dem Datenblatt und anderen Quellen verglichen.

Da die Querschnittsfläche einer stehenden Zugprobe deutlich kleiner ist, als bei den liegenden Zugproben, führt dies zu dem Phänomen, dass die zuvor gedruckte Schicht nicht direkt abkühlen kann, ehe dann die nächste Schicht gedruckt wird. Es ist gelungen stehende Zugproben zu drucken, indem die Druckgeschwindigkeit deutlich verringert und die Lüfterdrehzahl auf 60 \% reduziert ist. Jedoch sind die Proben beim Versand zu \textit{L.B. Bohle} in dem Päckchen zerstört. Daher wird sich im Rahmen dieser Arbeit auf den Vergleich zwischen den liegenden Zugproben hinsichtlich der mechanischen Kennwerte bezogen.
Im Slicer sind dann in x- und y-Richtung um 17 \% vergrößert und in z-Richtung sind sie 20 \% größer gedruckt. Damit ist sichergestellt, dass die Zugproben im gesinterten Zustand die DIN-Norm einhalten. Die fertig gesinterten Zugproben werden dann in eine Zugprüfmaschine an beide Seiten eingespannt und bis zur Zerstörung gezogen. Die benötigte Kraft, sowie die entstehende Längung bilden die Grundlage für das entstehende Spannungs-Dehnungs-Diagramm.
Das zugehörige Spannungs-Dehnungs-Diagramm ist in \autoref{ZugprobenS} dargestellt, anzumerken ist an dieser Stelle, dass nur \autoref{ZugprobenS} Zugproben enthält, die im Rahmen dieser Arbeit gedruckt sind. Die anderen drei Ergebnisse \autoref{ZugprobenGS}, \autoref{ZugprobenP} und \autoref{ZugprobenT} sind Zugproben, die in der Firma \textit{L.B. Bohle} gedruckt sind. Sie dienen nur als Vergleich.
 
\section{Vergleich der ermittelten Werte}
\label{sec:Vergleich}

\subsection{Vergleich der Dichte}

Die errechnete, durchschnittliche Dichte aus \autoref{Dichte- und Schwindung} beträgt 6,87 \(\text{g/cm}^3\). Verglichen mit dem technischen Datenblatt aus \autoref{Datenblatt1.4404} ist die Dichte dort mit 8 \(\text{g/cm}^3\) angegeben. Dies entspricht einer Abweichung von -14 \%. In \Autocite{Quarto.2021} wurde ebenfalls die Dichte an FDM-gedruckten Stahlteilen gemessen, hier lag sie im Durchschnitt bei knapp über 7,2 \(\text{g/cm}^3\). Dies entspricht einer Abweichung von -10 \%. Dieser Wert ist mit den Ergebnissen dieser Arbeit nur bedingt vergleichbar. In \autocite{Quarto.2021} wurde mit \textit{Ultrafuse} von \textit{BASF} gearbeitet. Das Material entspricht ebenfalls Edelstahl der Norm 1.4404. Die Abweichungen von Soll- und Istwerten werden mit \autoref{Abweichungen} errechnet.

\begin{equation}
  \text{Prozentuale Abweichung} = \left( \frac{\text{Soll-Dichte} - \text{Ist-Dichte}}{\text{Soll-Dichte}} \right) \times 100
  \label{Abweichungen}
\end{equation}
  


\subsection{Vergleich der Werte aus dem Zugversuch}

In \autoref{tab:ErgebnisseZugversuch} sind die relevantesten Ergebnisse aus den Zugversuchen aufgelistet. R$_{\textsubscript{p0.2}}$ stellt die Streckgrenze dar. R$_{\textsubscript{m}}$ steht für die Zugfestigkeit und A$_{\textsubscript{25 mm}}$ gibt die Bruchdehnung an. Verglichen werden hier die Werte für die Zugversuche der liegend gedruckten Proben. In der ersten Zeile sind die Proben, die im Rahmen dieser Arbeit gedruckt wurden aufgelistet \textit{(FH-Münster)}. In der zweiten Zeile sind die Proben der Firma \textit{L.B. Bohle} dargestellt. In der letzten Zeile \glqq Abweichungen\grqq sind die Abweichungen der ersten und zweiten Zeile von dem Datenblatt dargestellt.

\begin{table}[h]
  \centering
  \caption{Ergebnisse Zugversuch}
  \begin{tabular}{cccc}
  \toprule
  \textbf{Serie} & \textbf{R$_{\textsubscript{p0.2}}$ [MPa]} & \textbf{R$_{\textsubscript{m}}$ [MPa]} & \textbf{A$_{\textsubscript{25 mm}}$ [mm]} \\
  \midrule
  liegend (FH-MS) & 157 & 346 & 23,2 \\
  liegend (L.B. B.) & 178 & 400 & 25,4 \\
  Datenblatt 1.4404 \autocite{AGSTSteel:1.4404} & $\leq$ 200 & 500-700 & 40 \\
  \midrule
  Abweichungen (\%) 
    & -21,5 & -31,2 & -42,0 \\
    & -11,0 & -20,0 & -37,5 \\
  \bottomrule
  \end{tabular}
\label{tab:ErgebnisseZugversuch}
\end{table}


\chapter{Kapitel 5}

5.1 Zusammenfassung der Ergebnisse\\
5.2 Beantwortung der Forschungsfragen\\
5.3 Limitationen der Studie\\

\section{Ausblick}

Zur Verbesserung der Dichte kann der Durchmesser der Druckerdüse variiert werden. In \autoref{Quarto.2021} wurde eine 0,6mm Düse verwendet. Ebendfalls kann die Schrumpfung auch an größeren Würfeln gemessen werden. Dies lässt Rückschlüsse auf das Verhalten des Polymers an größeren Volumina zu.

5.4 Ausblick und zukünftige Arbeiten

%% ...
%%\chapter{Diskursiver Teil}
\label{Zusammenfassung}



%% Beginn vom Anhang ------------------------------------------------------------------------------
\appendix

%% Literaturverzeichnis einbinden. Die Quellen sind in der Datei "LiteraturDB.bib" gespeichert. ---
%% Das Literaturverzeichnis muss vorab über über Tools-->Befehle-->Biber erzeugt werden.--
\begin{flushleft}  % Schaltet den Blocksatz für das Literaturverzeichnis ab.
     \printbibliography
\end{flushleft}
%% Anhang einbinden -------------------------------------------------------------------------------
\chapter{Anhang}
\label{Anhang}
\begin{table}[htbp]
    \centering
    \caption{Parametertabelle Prusa MK3S+\newline verwendet für Druck der Würfel und liegenden Zugproben}
      \begin{tabular}{llr}
      \toprule
      \textbf{Parameter} & \textbf{Wert} & \multicolumn{1}{l}{\textbf{Beschreibung}} \\
      \midrule
      Schichthöhe [mm] & \multicolumn{1}{r}{0,1} &  \\
      Drucktemperatur [°C] & \multicolumn{1}{r}{130} &  \\
      Druckgeschwindigkeit [mm/s] & \multicolumn{1}{r}{50} & \multicolumn{1}{p{12.555em}}{Liefert optisch gute Ergebnisse} \\
      Heizbetttemperatur [°C] & \multicolumn{1}{r}{40} &  \\
      Rückzugslänge [mm] & \multicolumn{1}{r}{0,6} & \multicolumn{1}{p{12.555em}}{Standardwert Cura} \\
      Materialflussrate [\%] & \multicolumn{1}{r}{100} &  \\
      Druckoberfläche & Flachkreppband & \multicolumn{1}{p{12.555em}}{Gute Ablöseeigenschaften mit Spachtel} \\
      Prozentsatz Aussenhaut überlappen [\%/mm] & 10/0,04 &  \\
      \bottomrule
      \end{tabular}%
    \label{Würfelparameter}%
  \end{table}


% Table generated by Excel2LaTeX from sheet 'Tabelle1'
\begin{table}[htbp]
    \centering
    \caption{Abmaße, Volumen, Gewicht und Dichte des Würfels als Grünteil}
      \begin{tabular}{lrrrrrrr}
      \toprule
      \textbf{Würfel:} & \multicolumn{1}{l}{\textbf{Z (Höhe) [mm]:}} & \multicolumn{1}{l}{\textbf{Länge1 [mm]:}} & \multicolumn{1}{l}{\textbf{Länge2 [mm]:}} & \multicolumn{1}{l}{\textbf{Gesamtvolumen [mm3]:}} & \multicolumn{1}{l}{\textbf{Gesamtvolumen [cm3]:}} & \multicolumn{1}{l}{\textbf{Gewicht [g]:}} & \multicolumn{1}{l}{\textbf{Dichte [g/cm3]}} \\
      \multicolumn{1}{r}{1} & 10,03 & 10,13 & 10,2  & 1036,36 & 1,0364 & 4,775 & 4,607 \\
      \multicolumn{1}{r}{1} & 10,01 & 10,12 & 10,2  & 1033,27 & 1,0333 & 4,765 & 4,612 \\
      \multicolumn{1}{r}{1} & 10,02 & 10,15 & 10,17 & 1034,32 & 1,0343 & 4,76  & 4,602 \\
      \multicolumn{1}{r}{1} & 10,04 & 10,14 & 10,19 & 1037,40 & 1,0374 & 4,761 & 4,589 \\
      \multicolumn{1}{r}{1} & 10,03 & 10,15 & 10,19 & 1037,39 & 1,0374 & 4,76  & 4,588 \\
      \multicolumn{1}{r}{1} & 10,05 & 10,15 & 10,14 & 1034,36 & 1,0344 & 4,768 & 4,610 \\
      \multicolumn{1}{r}{1} & 10,06 & 10,13 & 10,19 & 1038,44 & 1,0384 & 4,751 & 4,575 \\
      \multicolumn{1}{r}{1} & 10,03 & 10,16 & 10,16 & 1035,35 & 1,0354 & 4,763 & 4,600 \\
      \multicolumn{1}{r}{1} & 10,02 & 10,17 & 10,16 & 1035,34 & 1,0353 & 4,762 & 4,599 \\
      \textbf{Schnitt1:} & \textbf{10,03} & \textbf{10,14} & \textbf{10,18} & \textbf{1035,80} & \textbf{1,0358} & \textbf{4,7628} & \textbf{4,5982} \\
      \midrule
      \multicolumn{1}{r}{2} & 9,94  & 10,22 & 10,2  & 1036,19 & 1,0362 & 4,724 & 4,559 \\
      \multicolumn{1}{r}{2} & 10,01 & 10,2  & 10,22 & 1043,48 & 1,0435 & 4,731 & 4,534 \\
      \multicolumn{1}{r}{2} & 9,97  & 10,18 & 10,22 & 1037,27 & 1,0373 & 4,754 & 4,583 \\
      \multicolumn{1}{r}{2} & 9,93  & 10,16 & 10,18 & 1027,05 & 1,0270 & 4,741 & 4,616 \\
      \multicolumn{1}{r}{2} & 9,97  & 10,21 & 10,22 & 1040,33 & 1,0403 & 4,751 & 4,567 \\
      \multicolumn{1}{r}{2} & 9,95  & 10,21 & 10,22 & 1038,24 & 1,0382 & 4,737 & 4,563 \\
      \multicolumn{1}{r}{2} & 9,95  & 10,15 & 10,23 & 1033,15 & 1,0332 & 4,741 & 4,589 \\
      \multicolumn{1}{r}{2} & 9,94  & 10,21 & 10,17 & 1032,13 & 1,0321 & 4,738 & 4,591 \\
      \textbf{Schnitt2:} & \textbf{9,96} & \textbf{10,19} & \textbf{10,21} & \textbf{1035,98} & \textbf{1,0360} & \textbf{4,7396} & \textbf{4,5751} \\
      \bottomrule
      \end{tabular}%
    \label{Grünteilmaße}%
  \end{table}%
  

\end{document}
